\subsection{A validated COMSOL state-to-KAM learning certificate}
\label{sec:validated-integral-learning}

We add a controlled identification experiment for the same complete
two-mode magnetic Hamiltonian. COMSOL independently integrates its four
state equations for $64$ initial conditions, with $17$ state observations
on each dimensionless interval $[0,T]$, $T\simeq0.04$.
This experiment is explicitly a two-mode modal ODE experiment; the
earlier $64$-DOF finite-element trajectories remain a separate validation.
Only positions and velocities enter the new training labels.
Neither acceleration reconstruction nor a solver tolerance is used as
an error certificate.

For each observed orbit, fit a degree-eight polynomial curve
$c_\ell:[0,T]\to\mathbb R^4$ to the observed states.
Its coefficients are frozen as exact dyadic rationals, with its constant
and linear coefficients corrected exactly so that both endpoints equal
the exported observations $z_\ell^-,z_\ell^+$.
This construction uses no evaluation of the physical acceleration.
Let $Lz=X_{h_2}(z)$ and let
$V(z)=[J\nabla F_0(z),\ldots,J\nabla F_{34}(z)]$
be the nonlinear feature field of the same learned family. Define
\begin{equation}
 \begin{aligned}
 D_\ell&=\frac1T\int_0^T V(c_\ell(\tau))\,d\tau,\\
 y_\ell&=\frac{z_\ell^+-z_\ell^-}{T}
          -\frac1T\int_0^T Lc_\ell(\tau)\,d\tau .
 \end{aligned}
 \label{eq:validated-integral-observations}
\end{equation}
The integral design and labels are therefore functions of the observed
state curves, not instantaneous force-balance labels.

\begin{proposition}[Integral loss, dynamic defect and coefficient error]
\label{prop:integral-transfer}
Let $\eta_\star=(s_4,0,\ldots,0)$ represent the declared full-cosine target,
and let $m$ be the number of observed curves. Put
\begin{equation}
 \begin{aligned}
 e_\ell&=\frac1T\int_0^T
       \bigl(\dot c_\ell-X_{h_\star}(c_\ell)\bigr)\,d\tau,\\
 G_I&=\frac1{4m}\sum_{\ell=1}^m D_\ell^\top D_\ell,\\
 r_I^2&=\frac1{4m}\sum_{\ell=1}^m
                   \|D_\ell\widehat\eta-y_\ell\|_2^2 .
 \end{aligned}
 \label{eq:integral-defect-gram}
\end{equation}
If $\lambda_{\min}(G_I)\geq\underline\lambda>0$ and
$(4m)^{-1}\sum_\ell\|e_\ell\|_2^2\leq\epsilon_I^2$, then
\begin{equation}
 \|\widehat\eta-\eta_\star\|_2
 \leq b_I:=\frac{r_I+\epsilon_I}{\sqrt{\underline\lambda}} .
 \label{eq:integral-coefficient-transfer}
\end{equation}
The curves need not be exact solutions of the physical differential
equation. Their explicitly verified dynamic defect supplies the required
error budget.
\end{proposition}

\begin{proof}
Exact endpoint matching and the fundamental theorem of calculus give
$y_\ell-D_\ell\eta_\star=e_\ell$.
The triangle inequality in the stacked, $1/(4m)$-normalized Euclidean
norm gives
$\|D(\widehat\eta-\eta_\star)\|_2/\sqrt{4m}\leq r_I+\epsilon_I$.
The Gram lower bound proves~\eqref{eq:integral-coefficient-transfer}.
\end{proof}

The defect is independently evaluated against the explicitly declared
physical target. Thus the proposition does not infer arbitrary unknown
model error from a small training loss. It certifies consistency with
this target using additional, verified physical information.

\paragraph{Validated integrals and numerical error.}
All polynomial feature integrals are evaluated by exact polynomial
integration with Arb ball coefficients. For $\vartheta=d^\top(c_\ell)_X$,
the magnetic feature uses the odd Taylor terms of
$\sin\vartheta-\vartheta$ through degree $19$, together with
\begin{equation}
 \begin{aligned}
 S_{19}(\vartheta)&=
 \sum_{\substack{3\leq j\leq19\\j\ {\rm odd}}}
       \frac{(-1)^{(j-1)/2}\vartheta^j}{j!},\\
 |\sin\vartheta-\vartheta-S_{19}(\vartheta)|
 &\leq\frac{R^{21}\cosh R}{21!},\quad |\vartheta|\leq R .
 \end{aligned}
 \label{eq:integral-sine-remainder}
\end{equation}
The computed maximum remainder is below $4.8752\times10^{-43}$.
The remainder is included in every affected feature integral and dynamic
defect; the certified function remains the complete cosine model.
The exact identity $y-D\eta_\star=e$ is also checked componentwise by
independent ball evaluations.

The interval Gram inverse is computed at $256$-bit precision.
Since the exact Gram matrix is real symmetric and positive semidefinite,
a validated inverse makes it positive definite and gives
$\underline\lambda=1/\|G_I^{-1}\|_\infty$.
Decoder rounding is retained when the real readout coefficients are
exported as a cosine multiplier and a quartic polynomial.
With $B=[q_\star/s_4\ U]$, this gives
\begin{equation}
 \begin{aligned}
 \|A_{\widehat\eta}-A_\star\|_2
 &\leq\sqrt{19}\bigl(\|B\|_2b_I+\delta_{\rm dec}\bigr),\\
 \|h_{\widehat\eta}-h_\star\|_r
 &\leq b_I B_F(r)+\delta_{h,\rm dec}(r).
 \end{aligned}
 \label{eq:integral-analytic-transfer}
\end{equation}
Here $\|B\|_2\leq\|B^\top B\|_\infty^{1/2}$ is verified from the
exported feature matrix; exact numerical orthogonality is unnecessary.
The analytic bound controls every derivative on a smaller polydisc by
Cauchy's estimate:
\[
 \begin{aligned}
 &\|\partial^\alpha(h_{\widehat\eta}-h_\star)\|_{r'}\\
 &\quad\leq\alpha!(r-r')^{-|\alpha|}
       \|h_{\widehat\eta}-h_\star\|_r .
 \end{aligned}
\]

\begin{table}[t]
\centering
\caption{Outward bounds for the independent COMSOL integral-learning
experiment. The MSE has a strictly positive certified lower bound.}
\label{tab:integral-certificate}
\begin{tabular}{lr}
\hline
Quantity & Certified bound\\
\hline
Training MSE & $[3.6030,3.6031]\,10^{-23}$\\
$r_I$ & $\leq6.0026\,10^{-12}$\\
$\epsilon_I$ & $\leq6.8835\,10^{-12}$\\
$\lambda_{\min}(G_I)$ & $\geq2.9302\,10^{-7}$\\
$\|\widehat\eta-\eta_\star\|_2$ & $\leq2.3805\,10^{-8}$\\
Twist error & $\leq1.0377\,10^{-7}$\\
$\sigma_{\min}(A_\star)$ & $\geq4.9032\,10^{-4}$\\
$\|h_{\widehat\eta}-h_\star\|_{0.3}$ & $\leq1.3310\,10^{-9}$\\
\hline
\end{tabular}
\end{table}

The target is unchanged. The singular-value lower bound in
Table~\ref{tab:integral-certificate} is a sharper bound for the same
two-by-two symmetric matrix:
$\sigma_{\min}(A_\star)\geq|\det A_\star|/\|A_\star\|_\infty$.
The previously reported, smaller lower bound remains valid.
The strict twist margin, fixed quadratic part and complete analytic
tail therefore give local positive-measure KAM families for this actual
fitted Hamiltonian by the normal-form argument above.

\paragraph{An independent bound on trajectory error.}
The observation curves deviate from every exported sample by at most
$5.2538\times10^{-14}$ and satisfy
\[
 \sup_\tau\|\dot c_\ell-X_{h_\star}(c_\ell)\|_2
 \leq1.0973\times10^{-8}.
\]
The linear modal field $L$ is skew symmetric in these balanced
coordinates. On the real box $|z|_\infty<0.3$, the nonlinear part has
Euclidean Lipschitz constant
\[
 \|d\|_2^2\frac{(0.3\|d\|_1)^2}{2}
 <0.002084<0.01 .
\]
For the exact physical solution $z_\ell$ with
$z_\ell(0)=c_\ell(0)$, variation of constants and Gronwall's inequality
therefore give
\[
 \begin{aligned}
 &\sup_{0\leq\tau\leq T}\|z_\ell(\tau)-c_\ell(\tau)\|_2\\
 &\quad\leq T e^{0.01T}
       \|\dot c_\ell-X_{h_\star}(c_\ell)\|_{\infty,2}\\
 &\quad<4.3907\times10^{-10}.
 \end{aligned}
\]
The verified curve bound $|c_\ell|_\infty<0.2000$ closes the
bootstrap inside the radius-$0.3$ box. This is an a posteriori bound
on the independently observed trajectories, not an interpretation of
COMSOL's requested time-stepping tolerances.

\paragraph{The same fitted function at millimetre amplitude.}
The frozen network has cosine multiplier
$a=0.9999996888354415$.
Its complete exported coefficients are also used in a separate,
$512$-bit a posteriori torus verification with the frequency in
Eqs.~\eqref{eq:certified-invariance}--\eqref{eq:certified-quadratic-frequency}
and the theorem constants of Theorem~\ref{thm:beam-finite-tori}.
The verified analytic invariance residual is at most
$2.7310\times10^{-77}$, the KAM smallness quotient is at most
$0.031724$, and the inverse averaged torsion is below $80526$.
The invariant torus lies within $2.0894\times10^{-41}$ of its exact
Fourier candidate. The tip-displacement witness is at least
$1.8977904$ mm, and the whole reconstructed centerline is bounded by
$1.8977962$ mm for every phase and all time.
Thus the nonzero integral-loss certificate and the finite-amplitude
KAM certificate concern exactly the same independently trained function.
The first proves coefficient, analytic and twist control; the second
checks the additional invariance-residual condition needed at the stated
finite amplitude.

The full-FE trajectory fit has an independent a posteriori torus
certificate, but its integral observation matrix yields a coefficient
bound above $0.025$ before accounting for continuous-quadrature and
model errors. This bound does not satisfy the coefficient condition
used here. The independently observed two-mode trajectories provide
the defect bounds needed for the complete integral-learning theorem;
model-reduction errors remain part of the separate full-FE evaluation.
